\begin{proposition}[C.3.17 (Proved): a marked circle extension over the cusp]
Let $V$ be a sufficiently small closed cusp body and
$\overline V=V/S^1_\delta$. Its smooth orbit structure can be
chosen compatibly with C.3.14 so that there is a diffeomorphism
\[
 \overline V\simeq E_{I_1}\times S^1,
\]
where $E_{I_1}$ is a closed neighbourhood of one nodal fibre in
an elliptic surface. More precisely, the actual positive
$\gamma$-translation circle action on the outer cusp collar
extends to a smooth circle action $\rho$ on $\overline V$.
The extension is free. It is not asserted to preserve the
fixed sphere $F\subset\overline V$ or the original interior
map to the $q$-disc.

Consequently the actual annular map $\psi$ of C.3.10 is the
endpoint of an ambient isotopy of $Q$ whose restriction to $F$
is the loop of embedded spheres $t\mapsto\rho_t(F)$.
After the marked deformation used in the proof, its trace is
the normalization immersion of a nodal elliptic curve times a
circle: it has one transverse double circle and two distinct
preimage circles. The regular cusp three-torus with period
basis $(\gamma,u,w)$ represents a primitive generator of
$H_3(\overline V;\Z)$; in particular it does not bound a
four-chain inside the cusp body.
\end{proposition}
\begin{proof}
We first remove the bounded period units by a marked smooth
deformation. On the cusp write, as in C.3.5,
\[
 Z(q)=sB_0+C(q),\qquad
 B_0=\begin{pmatrix}0&1\\-1&0\end{pmatrix},\qquad
 q=e^{2\pi i s}.
\]
For $0\leq a\leq1$ replace $C(q)$ by $(1-a)C(q)$, keeping
the same $A_2$ triangulation and the real periods $w,\delta$.
This is a deformation of the actual toric period quotient,
not a replacement inferred from its homology. The two period
automorphisms still have the same integral translations of the
fan and commute; only their nowhere-zero holomorphic unit
multipliers change. They depend smoothly on $a$, and even
holomorphically on a small complex neighbourhood of this
interval. On a sufficiently small punctured $q$-disc the
imaginary period matrices are invertible uniformly in $a$:
their leading part is $(\operatorname{Im}s)B_0$, while the
remaining matrices are uniformly bounded.

The fixed toric charts, with this extra parameter, therefore
give a smooth family of the same Mumford cusp quotients.
There is no new stabilizer at the central fibre: a nonzero
period translates every height-one ray by a nonzero lattice
vector, and cannot preserve any finite cone of the fan.
The usual finite fundamental collection of toric charts gives
proper discontinuity, uniformly for the compact parameter
interval. The central family has the constant compact
normalization $dP_6$ with its smoothly varying opposite-side
identifications. Thus it has a proper representative over a
uniform smaller $q$-disc. One may see the uniformity directly
by taking a relatively compact neighbourhood of this compact
central family: its boundary is disjoint from $q=0$, and after
shrinking $|q|$ no fibre meets that boundary. The nonzero
fibres are connected compact tori, so the selected representative
contains the whole nearby fibre. Denote these compact smooth
cusp bodies by $V_a$.

Their $S^1_\delta$ actions form a smooth semifree family, with
a smooth family of fixed two-spheres. Choose an equivariant
Ehresmann transport with the following two prescribed local
forms. On the outer collar keep $q$ and all real coefficients
in the marked period basis fixed. This is a smooth transport
there and carries the actual $\gamma$-translation action to
the corresponding action for every $a$. Near the fixed spheres,
choose family equivariant normal tubular coordinates. After
conjugating the negative-weight line the normal representation
is scalar multiplication on $\mathbb C^2$. Choose Hermitian
metrics and a unitary connection, and make the transport
linear and unitary in these normal coordinates. The collar
and fixed tubes are disjoint. A circle-invariant partition of
unity extends and patches the prescribed horizontal fields;
each has derivative one on $a$, so their patched field does
also. Compactness gives its flow for the entire interval.

This particular transport induces a smooth transport of orbit
spaces. On the free locus this follows from the principal-bundle
quotient. In the fixed tubes, unitary normal transport induces
the usual smooth linear rotation in the Hopf $\mathbb R^3$
coordinates. Thus it is smooth there as well. This explicit
choice avoids any assertion that every equivariant smooth map
has a smooth map in an arbitrarily preassigned quotient atlas.
The Hopf smooth structures and tubular choices are legitimate
ones for the orbit construction in C.3.14; the annular map is
supported in the free collar, where these structures coincide.
We obtain a marked orbit diffeomorphism
\[
 (\overline V,\text{boundary }\gamma\text{-action})
 \simeq(\overline V_1,\text{boundary }\gamma\text{-action}).
\]
The transport is not required to preserve $q$ in the interior.

At $a=1$ the period matrix is $sB_0$. The complex torus is
the product of two Tate elliptic curves, and the square toric
model is the fibre product of their nodal families. The chosen
$A_2$ triangulation is its chosen small resolution. Project
this resolution to the first Tate family $E_{I_1}$, and let
$S^1_\delta$ act on the second factor. The projection descends
to a proper smooth submersion
\[
 p:\overline V_1\longrightarrow E_{I_1}
\]
with oriented circle fibres. Here is the local check at the
only place where this assertion is not immediate. Write the
fibre product as $xy=uv$, with $q=xy=uv$, and take the small
resolution chart
\[
 u=xA,\qquad y=Av,
\]
whose coordinates are $(x,A,v)$. The weights of $(A,v)$ are
$(+1,-1)$ and $x$ has weight zero. Its Hopf quotient coordinates
are
\[
 \bigl(x,y=Av,d=(|A|^2-|v|^2)/2\bigr),
\]
and the descended projection is simply $(x,y,d)\mapsto(x,y)$.
The second chart, $x=uB$, $v=By$, gives the identical check
with $d'=(|u|^2-|B|^2)/2$. The local orbit fibre at $x=y=0$
consists of the two orbit rays joined by the quotient interval
of the exceptional projective line, and is a smooth real line.
Globally it closes to the second Tate orbit circle. At fixed
points away from this chart the model is likewise
$(p,A,v)\mapsto(p,Av)$, whose quotient is the same smooth
projection. We make the compatibility of these quotient charts
explicit. On overlaps near the same fixed branch the two weight
coordinates have the form
\[
 A_j=a_{ij}(p,z)A_i,\qquad v_j=b_{ij}(p,z)v_i,
 \qquad z=A_iv_i,
\]
where the units and the change of the tangential coordinate
$p$ are pulled back from the projection base. The two fixed
lifts over the node have disjoint fixed tubular neighbourhoods;
they require no overlap identification at the node itself.
Choose Hermitian metrics for the two weight line bundles on
these base charts, retaining the separated lifts of the fixed
normalization. If their positive local coefficients are
$h_{A,i},h_{v,i}$, their transition rules are
\[
 h_{A,j}|a_{ij}|^2=h_{A,i},\qquad
 h_{v,j}|b_{ij}|^2=h_{v,i}.
\]
The weighted Hopf coordinate
\[
 d_i=\tfrac12\bigl(h_{A,i}|A_i|^2-h_{v,i}|v_i|^2\bigr)
\]
therefore agrees with $d_j$ on overlaps. These are legitimate
Hopf coordinates: rescale the two normal variables by
$\sqrt{h_{A,i}}$ and $\sqrt{h_{v,i}}$. Their product is
$\sqrt{h_{A,i}h_{v,i}}\,z$, a smooth local change of the
complex coordinate $z$ with invertible derivative at $z=0$.
Thus $(p,z,d_i)$ are smooth quotient coordinates, their
overlap maps are smooth, and the projection remains
$(p,z,d_i)\mapsto(p,z)$. This specifies the Hopf smoothing
used for the transport above, without a general smoothness
claim for arbitrary equivariant quotient maps.

The section of the second Tate family represented by $z_2=1$
avoids its node. It therefore gives a section of the fibre
product which lifts through the small resolution, and then
gives a section of $p$. Hence this oriented circle bundle is
trivial. This proves the asserted product diffeomorphism by
an actual projection and section.

We can match the boundary action exactly, not just its homology
class. On the boundary of this circle bundle the positive
$\gamma$ action is the circle translation in the second Tate
orbit fibre. Its invariant vertical length form has fibre
integral one. Extend that length form to any smooth positive
vertical length form on the entire oriented circle bundle,
using a collar cutoff. Divide on each fibre by its integral;
this is a smooth positive function of the base because $p$
is a proper submersion. The dual vertical vector field now has
period exactly one, and agrees with the specified vector field
on the boundary collar. Its flow defines the desired free
circle action. Transport it back to $\overline V$ by the marked
orbit diffeomorphism.

Choose an invariant function $\chi$ equal to one in the inner
body and zero outside the outer collar, with the same collar
profile as in C.3.10. Such a function can be chosen because
the action agrees with the actual $q$-fibre translations there;
extend it constantly by one on the remaining interior. The
maps
\[
 H_t(z)=\rho_{t\chi(z)}(z),\qquad 0\leq t\leq1,
\]
are diffeomorphisms, since $\chi$ is invariant. They extend
by the identity to $Q$. Their endpoint is the actual annular
map: on the inner body it is $\rho_1=\id$, and on the collar
it is translation by $\chi\gamma$. On the fixed sphere
$\chi=1$, so this isotopy traces precisely $\rho_t(F)$.
It is not an isotopy relative to the fixed sphere.

In the zero-unit model the fixed sphere projects to the
normalization of the first nodal elliptic curve. Over a smooth
point of that curve it meets the orbit circle once, while the
two preimages of its node give two distinct points of the
orbit circle over the node. Its circle sweep is therefore the
normalization immersion of the nodal curve times the circle.
The local form $(x,y,d)$ with $q=xy$ gives its two transverse
branches at the double circle. Finally $E_{I_1}$ retracts onto
its nodal elliptic curve, whose fundamental two-cycle is the
degree-one specialization of a regular elliptic fibre. The
K\"unneth theorem now gives
\[
 H_3(E_{I_1}\times S^1;\Z)
 =H_2(E_{I_1};\Z)\otimes H_1(S^1;\Z)=\Z,
\]
generated by that regular fibre times the marked circle.
Boundary marking transports this three-torus to the actual
$(\gamma,u,w)$ torus, proving the last assertion.
\end{proof}

\begin{remark}[Scope]
The raw fractional period translations need not extend across
the central fibre. The action above modifies their interior
orbits. Its loop moves the fixed sphere, which is exactly the
feature retained by the annular relative mapping class. Thus
the construction does not trivialize that class or evaluate
$J_\psi$. In particular the nonbounding three-torus prevents
discarding its stable framing contribution by a chain lying
entirely in the cusp body.
\end{remark}
